\documentclass[a4paper,12pt]{article}
\usepackage[slovene]{babel}
\usepackage[utf8]{inputenc}
\usepackage[T1]{fontenc}
\usepackage{lmodern}

\usepackage{amsmath}

\setlength{\parindent}{0mm}

\begin{document}

\section*{Matrike}
% Oblikujte determinanto matrike.
% Vsebina matrike je že pripravljena v komentarju spodaj.
% V pomoč naj vam bo Overleafova dokumentacija o matrikah:
% https://www.overleaf.com/learn/latex/Matrices
Izračunaj determinanto
\[
\begin{vmatrix}
    -1 & 4 & 4 & -2 \\
    1 & 4 & 5 & -1 \\
    1 & 4 & -2 & 2 \\
    3 & 8 & 4 & 3 \\
\end{vmatrix}
\]

\section*{Okolje \texttt{align}}
% Okolje align je namenjeno poravnavi enačb.
% Če ne želite, da se enačbe oštevilčijo, uporabite okolje align*.
% Oblikujte spodnje enačbe z okoljem align*.
% Če naredite po en kurzor na začetku vsake vrstice,
% boste lahko oblikovali vse tri vrstice hkrati.
Nariši grafe funkcij:

\begin{align*}
y &= x^2 - 3|x| + 2 & y &= 3 \sin(\pi+x) - 2\\
y &= \log_2(x-2) + 3 & y &= 2 \sqrt{x^2+15} + 6\\
y &= 2^{x-3} + 1 & y &= \cos(x-3) + \sin^2(x+1)\\
\end{align*}


\section*{Okolje \texttt{multline}}
% Oblikujte spodnjo enačbo z okoljem multline*,
% da bo oblikovana kot v rešitvi.
Poišči vse rešitve enačbe
\begin{multline*}
   (1+x+x^2) \cdot (1+x+x^2+x^3+\ldots+x^9+x^{10}) =\\
=(1+x+x^2+x^3+x^4+x^5+x^6)^2.
\end{multline*}

\section*{Okolje \texttt{cases}}
% Oblikujte spodnji funkcijski predpis z okoljem cases,
% da bo oblikovan kot v rešitvi.
Dana je funkcija
\[
f(x) = \begin{cases}
   \frac{3x^2y-y^3}{x^2+y^2}; &(x,y) \neq (0,0)\\
   a; &(x,y) = (0,0)\\
\end{cases}
\]
\begin{itemize}
\item Določi \(a\), tako da izračunaš limito \( \lim_{(x,y)\to(0,0)} f(x,y) \).
\item Izračunaj parcialna odvoda \(f_x(x,y)\) in \(f_y(x,y)\).
\end{itemize}

\end{document}
